\chapter*{Glossário}
\addcontentsline{toc}{chapter}{Glossário}
\label{glossario}

Este glossário reúne, em ordem alfabética, os termos técnicos de estatística, aprendizado de máquina e medicina usados ao longo do trabalho, com definições curtas. Os conceitos que exigem explicação mais detalhada são desenvolvidos no Capítulo~\ref{cap:fundamentacao}.

\begin{longtable}{@{}p{3.8cm}p{10.7cm}@{}}
AED (Análise Exploratória de Dados) & Conjunto de técnicas descritivas, como resumos numéricos, gráficos e testes estatísticos, usado para conhecer a estrutura e a qualidade dos dados antes da modelagem. \\[0.15cm]
Albumina & Principal proteína produzida pelo fígado e presente no sangue; seu nível baixo sugere comprometimento do órgão. \\[0.15cm]
ALT/TGP e AST/TGO & Enzimas presentes nas células do fígado que passam ao sangue quando elas são lesadas; níveis altos sugerem lesão hepática. \\[0.15cm]
Amostra sintética & Observação criada artificialmente, como no SMOTE, e não coletada de um paciente real. \\[0.15cm]
Árvore de decisão & Modelo que classifica por meio de uma sequência de perguntas sobre as variáveis, como um fluxograma; cada pergunta é um nó, e cada resultado final é uma folha. \\[0.15cm]
Assimetria (\textit{skewness}) & Grau em que a distribuição de uma variável deixa de ser simétrica em torno do centro; na assimetria positiva, há uma cauda longa de valores altos. \\[0.15cm]
AUC-ROC & Área sob a curva ROC; probabilidade de que um paciente doente sorteado ao acaso receba do modelo uma probabilidade maior que a de um saudável. Vale 0,5 para um classificador ao acaso e 1 para um perfeito. \\[0.15cm]
\textit{Bagging} & Estratégia de \textit{ensemble} em que vários modelos são treinados de forma independente, cada um com uma amostra diferente dos dados, e suas previsões são combinadas por votação. \\[0.15cm]
\textit{Baseline} (linha de base) & Modelo simples usado como referência para verificar se modelos mais complexos trazem ganho. \\[0.15cm]
\textit{Beeswarm} e \textit{waterfall} & Gráficos de resumo do SHAP: o primeiro mostra um ponto por paciente em cada variável; o segundo decompõe a previsão de um único paciente a partir do valor base. \\[0.15cm]
Bilirrubina & Pigmento amarelo resultante da degradação da hemoglobina, eliminado pelo fígado; níveis altos sugerem disfunção hepática ou obstrução biliar. \\[0.15cm]
Bonferroni (correção de) & Ajuste dos valores-$p$ quando vários testes são feitos, que multiplica cada valor-$p$ pelo número de testes para limitar os falsos positivos. \\[0.15cm]
\textit{Boosting} & Estratégia de \textit{ensemble} em que os modelos são treinados em sequência e cada um procura corrigir os erros dos anteriores. \\[0.15cm]
\textit{Bootstrap} & Amostra sorteada com reposição do conjunto de dados, com o mesmo tamanho dele; alguns registros aparecem repetidos e outros ficam de fora. \\[0.15cm]
\textit{Boxplot} & Gráfico que resume a distribuição de uma variável: a caixa vai do primeiro ao terceiro quartil, a linha interna marca a mediana, os bigodes se estendem até os valores não atípicos e os pontos isolados são valores extremos. \\[0.15cm]
Busca aleatória (\textit{Random Search}) & Otimização de hiperparâmetros que sorteia um número fixo de combinações dentro de um espaço de busca. \\[0.15cm]
Caixa-preta (\textit{black box}) & Modelo cujo funcionamento interno é complexo demais para ser inspecionado diretamente. \\[0.15cm]
Calibração & Grau em que as probabilidades previstas correspondem às frequências reais; um modelo bem calibrado que prevê 30\% de chance acerta o evento em cerca de 30\% dos casos. \\[0.15cm]
Chance (\textit{odds}) & Razão entre a probabilidade de um evento ocorrer e a de não ocorrer; probabilidade de 0,75 corresponde a chance de 3 para 1. \\[0.15cm]
Classificador trivial (ingênuo) & Classificador que ignora os preditores e prevê sempre a classe majoritária. \\[0.15cm]
Correlação de Pearson e de Spearman & A de Pearson mede o quanto os pontos de duas variáveis se aproximam de uma reta (varia de $-1$ a $+1$); a de Spearman faz o mesmo sobre as posições dos valores quando ordenados, sendo menos sensível a valores extremos. \\[0.15cm]
Curva ROC & Gráfico da sensibilidade em função da taxa de falsos positivos, para todos os limiares de decisão possíveis. \\[0.15cm]
DeLong (teste de) & Método estatístico que estima a incerteza da AUC-ROC e compara a AUC-ROC de dois modelos avaliados nos mesmos pacientes. \\[0.15cm]
Desbalanceamento & Diferença grande entre os números de observações das classes. \\[0.15cm]
Efeito Rashomon & Existência de vários modelos com desempenho praticamente igual que explicam os mesmos dados de formas diferentes, por exemplo atribuindo importâncias distintas às variáveis. \\[0.15cm]
\textit{Ensemble} & Conjunto de modelos cujas previsões são combinadas em uma única previsão. \\[0.15cm]
Especificidade & Proporção de pacientes saudáveis corretamente identificados como saudáveis. \\[0.15cm]
Estratificação & Divisão dos dados que mantém, em cada parte, a mesma proporção de classes da base original. \\[0.15cm]
Falso negativo (FN) & Paciente doente classificado como saudável. \\[0.15cm]
Falso positivo (FP) & Paciente saudável classificado como doente; em testes de hipótese, rejeitar uma hipótese nula verdadeira. \\[0.15cm]
\textit{Fold} & Cada uma das partes em que o conjunto de treino é dividido na validação cruzada. \\[0.15cm]
Folha & Nó final de uma árvore de decisão, que indica a classe ou o valor previsto. \\[0.15cm]
Fosfatase alcalina & Enzima associada às vias biliares; níveis altos sugerem obstrução biliar. \\[0.15cm]
Globulina & Grupo de proteínas do sangue; a razão albumina/globulina (A/G) compara a albumina com as demais proteínas e tende a diminuir na doença hepática. \\[0.15cm]
\textit{Gradient boosting} & Variante do \textit{boosting} em que cada nova árvore é ajustada na direção que mais reduz a função de perda. \\[0.15cm]
Gráfico de dependência (SHAP) & Gráfico do valor SHAP de uma variável em função do valor da própria variável. \\[0.15cm]
Hiperparâmetro & Configuração de um algoritmo definida antes do treinamento, que não é aprendida dos dados. \\[0.15cm]
Hipótese nula & Afirmação, geralmente de ausência de diferença ou de efeito, que um teste de hipótese procura contrariar. \\[0.15cm]
Importância global e local (SHAP) & A global é a média do efeito de uma variável sobre todos os pacientes; a local descreve seu efeito sobre um único paciente. \\[0.15cm]
Imputação & Preenchimento de valores ausentes com valores estimados. \\[0.15cm]
Interpretabilidade & Capacidade de entender por que um modelo produz determinada previsão. \\[0.15cm]
Intervalo de confiança & Faixa de valores plausíveis para uma quantidade estimada, dada a variabilidade da amostra. \\[0.15cm]
KDD (\textit{Knowledge Discovery in Databases}) & Processo de descoberta de conhecimento em bases de dados, que vai da seleção e preparação dos dados à interpretação dos padrões encontrados. \\[0.15cm]
\textit{Leaf-wise} & Forma de crescer uma árvore expandindo sempre a folha cuja divisão mais reduz a perda, e não nível por nível. \\[0.15cm]
Limiar clínico & Maior limiar que atinge uma sensibilidade mínima definida previamente (neste trabalho, 90\%). \\[0.15cm]
Limiar de decisão & Valor de probabilidade a partir do qual uma observação é classificada como positiva. \\[0.15cm]
Logaritmo da chance (\textit{log-odds}) & Logaritmo natural da chance; escala em que a Regressão Logística e o \textit{boosting} somam suas contribuições. \\[0.15cm]
Matriz de confusão & Tabela que cruza as classes previstas com as classes reais e contabiliza os acertos e os erros de cada tipo. \\[0.15cm]
Monte Carlo (simulação de) & Método que estima uma quantidade por meio de sorteios aleatórios repetidos. \\[0.15cm]
\texttt{mtry} & Número de variáveis sorteadas a cada pergunta de uma árvore do \textit{Random Forest}. \\[0.15cm]
Multicolinearidade & Situação em que duas ou mais variáveis preditoras são fortemente correlacionadas e carregam praticamente a mesma informação. \\[0.15cm]
\textit{Outlier} (valor extremo) & Observação muito afastada da maioria dos dados. \\[0.15cm]
\textit{Pipeline} & Sequência de etapas de processamento e modelagem executada em código, da leitura dos dados à avaliação dos modelos. \\[0.15cm]
Posto (\textit{rank}) & Posição de um valor quando os dados são ordenados do menor para o maior; empates recebem o posto médio. \\[0.15cm]
Precisão & Proporção de doentes entre os pacientes classificados como doentes. \\[0.15cm]
Prevalência & Proporção de observações de uma classe em relação ao total. \\[0.15cm]
Previsão fora da partição (\textit{out-of-fold}) & Previsão feita para um paciente por um modelo treinado sem ele, na validação cruzada. \\[0.15cm]
\textit{Random Forest} & Modelo de \textit{ensemble} que combina muitas árvores de decisão treinadas com \textit{bagging} e com sorteio de variáveis a cada pergunta. \\[0.15cm]
Razão de chances (\textit{odds ratio}) & Fator pelo qual a chance de um evento é multiplicada quando um preditor aumenta em uma unidade, na Regressão Logística. \\[0.15cm]
Regressão Logística & Modelo estatístico para variáveis-alvo binárias que escreve o logaritmo da chance como uma soma ponderada dos preditores. \\[0.15cm]
Regularização & Penalização da complexidade de um modelo durante o treinamento, para reduzir o sobreajuste. \\[0.15cm]
Semente aleatória & Valor inicial do gerador de números aleatórios que torna os sorteios reproduzíveis. \\[0.15cm]
Sensibilidade (\textit{recall}) & Proporção de pacientes doentes corretamente identificados como doentes. \\[0.15cm]
SHAP & Método que decompõe a previsão de um modelo em contribuições de cada variável, com base nos valores de Shapley. \\[0.15cm]
Shapley (valor de) & Divisão justa de um ganho conjunto entre participantes, calculada pela média de suas contribuições sobre todas as ordens possíveis de entrada. \\[0.15cm]
SMOTE & Técnica de \textit{oversampling} que cria observações sintéticas da classe minoritária em pontos intermediários entre observações vizinhas. \\[0.15cm]
SMOTE-NC & Variante do SMOTE para bases com variáveis categóricas, em que o valor da variável categórica da observação sintética é definido a partir das observações de origem. \\[0.15cm]
Sobreajuste (\textit{overfitting}) & Situação em que o modelo se ajusta tanto aos dados de treino, inclusive ao ruído, que generaliza mal para dados novos. \\[0.15cm]
Taxa de falsos negativos e de falsos positivos & Proporção de doentes classificados como saudáveis e proporção de saudáveis classificados como doentes. \\[0.15cm]
Teste de Wald & Teste que avalia se um coeficiente da Regressão Logística difere de zero, comparando-o com seu erro-padrão. \\[0.15cm]
Transformação logarítmica & Substituição de cada valor pelo seu logaritmo, que comprime a cauda de valores altos sem alterar a ordem das observações. \\[0.15cm]
\textit{TreeSHAP} & Algoritmo rápido e exato para calcular valores SHAP em modelos baseados em árvores. \\[0.15cm]
Validação cruzada & Procedimento que divide o conjunto de treino em partes e reveza qual delas é usada para avaliar o modelo treinado com as demais. \\[0.15cm]
Valor base (SHAP) & Saída média do modelo sobre um conjunto de referência, ponto de partida da decomposição SHAP. \\[0.15cm]
Valor-$p$ & Probabilidade de se obter um resultado tão extremo quanto o observado se a hipótese nula fosse verdadeira. \\[0.15cm]
Vazamento de dados (\textit{data leakage}) & Uso, durante o treinamento, de informação do conjunto de teste, o que torna a avaliação otimista demais. \\[0.15cm]
Viés algorítmico & Tendência de um modelo a produzir resultados sistematicamente piores para determinados grupos de pessoas. \\[0.15cm]
XGBoost e LightGBM & Algoritmos de \textit{gradient boosting} com árvores de decisão, otimizados para velocidade e desempenho. \\[0.15cm]
\textit{Z-score} (padronização) & Valor menos a média, dividido pelo desvio-padrão; deixa variáveis com média zero e desvio-padrão um. \\
\end{longtable}
